\begingroup
\renewcommand{\arraystretch}{1.45}
\fancyhead[L]{\small BESTIMMUNG DER ANZIEHUNG}
\Needspace{40\baselineskip}
\AttractionArticle{18}
Durch die hier erläuterte Methode lassen sich sogar \textit{unbestimmte} Integrale (die beim Wert der Variablen $=0$ beginnen) mit größter Eleganz angeben. Nimmt man nämlich an, dass $T''$ durch $m',\ n',\ T'$ ebenso bestimmt werde wie $T'$ durch $m.\ n,\ T$, und ebenso wiederum $T'''$ durch $m'',\ n'',\ T''$, und so fort, so werden auch für jeden bestimmten Wert von $T$ die Werte der Glieder der Folge $T,\ T',\ T'',\ T'''$ usw. sehr rasch gegen einen Grenzwert $\Theta$ konvergieren, und es wird gelten:
\[
\int\frac{\mathrm{d}T}{\surd(mm\cos T^2+nn\sin T^2)}=\frac{\Theta}{\mu}
\]
\[
\begin{aligned}&\int\frac{(\cos T^2-\sin T^2)\mathrm{d}T}{\surd(mm\cos T^2+nn\sin T^2)}\\[7pt]&=-\frac{\nu\Theta}{\mu}\\[7pt]&\quad+\frac{\begin{gathered}\lambda'\cos T\sin T'+2\lambda''\cos T'\sin T''\\+4\lambda'''\cos T'\sin T'''+\mathrm{usw.}\end{gathered}}{\lambda\lambda}\end{aligned}
\]
\par Doch möge es genügen, dies hier nebenbei angedeutet zu haben, da es für unseren Zweck nicht erforderlich ist.

\Needspace{40\baselineskip}
\AttractionArticle{19}
Setzen wir nun $m=\surd(G+G')$, $n=\surd(G+G'')$, so werden sich die Werte der Größen $P,\ Q$ leicht auf die transzendenten Größen $\mu,\ \nu$ zurückführen lassen. Da nämlich $P,\ Q$ die Werte der Integrale
\SourceLeaf{364}
\[
\begin{gathered}\int\frac{\cos T^2\mathrm{d}T}{2\pi(mm\cos T^2+nn\sin T^2)^{\frac32}},\\[7pt]\int\frac{\sin T^2\mathrm{d}T}{2\pi(mm\cos T^2+nn\sin T^2)^{\frac32}}\end{gathered}
\]
bei Erstreckung von $T=0$ bis $T=360^{0}$ sind, ist zunächst unmittelbar ersichtlich, dass
\[
mmP+nnQ=\frac{1}{\mu}\quad .\, .\, .\, .\, .\, .\, .\, .\, .\, .\, .\, .\quad[24]
\]
\Needspace{36\baselineskip}
\par Ferner gilt
\[
\begin{aligned}
&\frac{(\cos T^2-\sin T^2)\mathrm{d}T}{2\pi\surd(mm\cos T^2+nn\sin T^2)}\\[7pt]&\quad+\frac{(mm\cos T^2-nn\sin T^2)\mathrm{d}T}{2\pi(mm\cos T^2+nn\sin T^2)^{\frac32}}\\[6pt]
&\qquad=\frac{(mm\cos T^4-nn\sin T^4)\mathrm{d}T}{\pi(mm\cos T^2+nn\sin T^2)^{\frac32}}\\[6pt]
&\qquad=\mathrm{d}.\frac{\cos T\sin T}{\pi\surd(mm\cos T^2+nn\sin T^2)}
\end{aligned}
\]
\par Integriert man diese Gleichung von $T=0$ bis $T=360^{0}$, so ergibt sich
\[
-\frac{\nu}{\mu}+mmP-nnQ=0\quad .\, .\, .\, .\, .\, .\, .\, .\, .\quad[25]
\]
\par Aus der Verbindung der Gleichungen 24, 25 schließen wir schließlich:
\[
P=\frac{1+\nu}{2mm\mu},\qquad Q=\frac{1-\nu}{2nn\mu}
\]
\begin{center}\rule{0.22\textwidth}{0.4pt}\end{center}
\clearpage
\endgroup
